\section{Discussion}
\label{sec:discussion}

\subsection{Why Real-Time and Multimodal Are Harder}

The reproducibility crisis documented in Section~\ref{sec:open_problems} is not
uniformly distributed across the guardrail literature.
It is \emph{especially severe} in real-time and multimodal settings, and the
reasons are structural.

\paragraph{Latency eliminates the strongest tools}
The most accurate safety mechanisms---LLM-as-judge pipelines, multi-step
constitutional reasoning---are incompatible with streaming delivery.
A post-generation judge that adds 1--2 seconds of latency is acceptable in a batch
API; it is a critical regression in a chat interface where first-token latency
drives user experience.
The real-time constraint therefore forces practitioners toward weaker, faster
mechanisms (small classifiers, heuristic filters) whose safety guarantees are
correspondingly harder to establish.

\paragraph{Modality boundaries break text-trained alignment}
Safety alignment as currently practised is fundamentally text-centric.
RLHF, constitutional AI, and safety context distillation all operate on the
text representation of model behaviour.
Multimodal fine-tuning introduces a new input pathway that was not considered
during alignment training.
The visual encoder, audio encoder, and their projections into the language model's
embedding space are potential entry points for misaligned behaviour, and aligning
them requires either extending the alignment training to cover multimodal inputs
explicitly or adding modality-specific classifiers as external guardrails.

\paragraph{The evaluation gap is widest here}
The evaluation-protocol fragility documented in Section~\ref{subsec:reproducibility}
is compounded in real-time and multimodal settings.
Latency measurements depend on hardware, batching strategy, and network conditions,
none of which are standardised.
Multimodal attack success rates depend on image resolution, encoding format, and
the specific VLM version, all of which vary across evaluations.
The result is that published claims in these domains are the least reproducible
in the entire guardrail literature.

\subsection{Implications for Guardrail Design}
\label{subsec:implications}

The following design principles are derived exclusively from the five surviving verified claims and from the structural null findings (streaming gap, audio/video gap documented in Section~\ref{sec:background}); they do not rest on unverified quantitative claims from the broader corpus.

\paragraph{Session-level context is mandatory}
The confirmed finding that multi-turn conversational priming amplifies jailbreaks
by 29--71\%~\cite{russinovich2024greatescalation} (derived from verified finding~4~\cite{russinovich2024greatescalation,li2024llm}) implies that any guardrail
operating only on the current turn is structurally insufficient against a capable
adversary.
Guardrails must maintain and inspect session-level context, which complicates
stateless inference architectures and increases memory requirements.

\paragraph{Modality-specific classifiers as a bridge}
Until multimodal alignment training is sufficiently mature, modality-specific
classifiers (image-content filters, audio-speech classifiers) can serve as a
first-pass guardrail layer before content reaches the language model backbone
(motivated by the audio/video absence documented in Section~\ref{subsec:audiovideo} and the taxonomy cell analysis in Section~\ref{sec:background}).
This approach is architecturally compatible with existing deployment patterns and
does not require retraining the underlying VLM.
Its limitation is that it cannot detect cross-modal attacks where the harmfulness
only emerges from the combination of text and image.
The evidential basis for this principle differs from the session-context and evaluation-first principles: it derives from a structural absence (no audio/video safety mechanisms in the corpus) rather than from a positive verified finding.
The prescription --- invest in modality-specific classifiers --- follows from the problem definition, not from demonstrated effectiveness; effectiveness data does not yet exist.

\paragraph{Evaluation-first development}
The reproducibility crisis suggests that guardrail development should adopt an
evaluation-first methodology: define the evaluation protocol before implementing
the mechanism (derived from verified finding~5: evaluation-protocol choices shift reported ASR by up to 30 percentage points~\cite{li2024salad}).
This mirrors test-driven development (TDD) in software engineering and is
aligned in spirit with the pre-registration movement in machine learning, which
advocates committing to evaluation criteria before observing results to prevent
post-hoc selection of favourable conditions~\cite{pineau2021improving}.
The analogy is useful but structurally inexact in three ways that require explicit
adaptation.
\emph{First}, unlike software unit tests that pass or fail deterministically,
safety evaluations have continuous outputs (attack success rates, false-positive
rates) that must be interpreted relative to a pre-registered threshold, not
compared post-hoc against the best available baseline.
\emph{Second}, because adversaries are adaptive, the attack set itself must be
updated before each evaluation cycle to include attacks published after the
guardrail was designed, preventing the evaluation from becoming a fixed target
that the mechanism eventually overfits --- a disanalogy absent from software
testing, where the test suite is typically stable.
\emph{Third}, software unit tests are independent by design: a failing test
isolates a specific code path.
Safety evaluations are \emph{correlated across prompts}: an attack that succeeds
on prompt $A$ is correlated with success on semantically similar prompt $B$,
because both exploit the same vulnerability in the model's decision boundary.
This correlation structure means that a guardrail evaluated on $n$ attack prompts
provides less independent evidence than $n$ unit tests; effective coverage requires
deliberately diverse attack sets that span harm categories and attack methodologies,
not merely large $n$.
These three adaptations are required for the TDD analogy to function prescriptively
in the safety evaluation context.
Concretely: a team developing a streaming safety classifier should specify---before
writing any classifier code---the exact attack prompts, generation budget,
judge model version, harm categories, and acceptable false-positive rate at which
the classifier will be considered successful.
The ``test suite'' should span at least the four attack families documented in
Section~\ref{sec:multimodal} (typographic, adversarial perturbation, indirect
injection, text-channel) to avoid exploiting correlation between prompt types.
This pre-registration prevents the post-hoc selection of evaluation conditions
that favour the proposed mechanism, which is the most common route to the
irreproducible results documented in this review.

\subsection{Are Five Verified Claims Sufficient?}

A legitimate challenge to this review is whether five verified findings constitute
a sufficient evidentiary foundation for the open-problems agenda developed in
Section~\ref{sec:open_problems}.
We argue that sufficiency must be judged against the intended use of the review,
not against an absolute threshold.

The five verified findings are not randomly distributed: they include the
field's most fully documented architectural mechanism (Llama Guard's classification
framing), its most reproducible implementation technique (safety context
distillation), its most concrete arbitration protocol (dual reward model with a
specified threshold), its most robust attack finding (multi-turn priming
amplification), and a meta-finding about evaluation methodology variability.
Together they identify \emph{what the field can currently stand on} — a narrow but
loadbearing set of verified facts.

The open problems in Section~\ref{sec:open_problems} are each grounded in one of
two ways: either in the absence of a verified claim in an area where the literature
asserts results (RT1--RT3, FG1, MM1), or in a verified structural fact that implies
an unmet requirement (MM2 follows directly from the Crescendo finding; FT1 follows
from the threat model established by Qi et al., whose numerics are unverified but
whose threat model is logically sound).
No open problem in the catalogue rests solely on an unverified claim treated as
established fact.

The alternative view — that five verified claims indicate the field is too
immature for a definitive SLR — is also defensible.
A survey that confirms so little may be premature as a synthesis but is premature
in a useful direction: it documents the evidentiary vacuum precisely enough that
future work can be measured against it.
The contribution of this review is therefore the methodology and the null result
as much as any positive synthesis.

\subsection{Limitations of This Review}

This review is subject to several limitations.

First, and most self-referentially, this review's central argument is a
reproducibility critique, yet its own screening, quality assessment, and
claim-verification were all performed by a single author.
The structured self-review protocol (Section~\ref{subsec:verification_protocol})
mitigates anchoring bias through temporally separated passes but cannot substitute
for genuinely independent multi-reviewer screening.
The following procedural compensations were applied to partially offset the single-author constraint: all inclusion/exclusion criteria were operationalised in a pre-specified codebook before screening began; an audit trail was maintained recording the primary reason for each exclusion; and the claim verification protocol's three-pass structure with mandatory unanimous agreement functions as a within-author reliability check, since a claim cannot survive if the same author's judgment differs across temporally separated passes.
These compensations do not substitute for independent multi-reviewer screening but do provide partial structural protection against anchoring bias in both the corpus assembly and the claim verification phases.
This is not merely a procedural limitation: the five verified claims that constitute
the headline empirical finding survived scrutiny under one codebook applied by one
author; a second independent coder applying the same codebook to a random sample
might return a higher or lower count, and the direction is unknown.
The five verified claims should therefore be understood as a \emph{lower bound}
on the reproducible findings in the literature under this codebook, not as the
definitive count.
A future replication using two or more independent reviewers with formally
computed pairwise inter-rater agreement (e.g.\ Cohen's $\kappa$ or Krippendorff's
$\alpha$) across the full corpus would strengthen confidence in both the inclusion
decisions and the claim-verification outcomes, and is the highest-priority
direction for follow-on methodological work.
The three design principles in Section~\ref{subsec:implications} are accordingly grounded only in the five verified claims and structural null findings; readers should treat them as expert synthesis derived from a narrow but robustly verified evidential base, not as conclusions from the full unverified corpus.

Second, the structured self-review protocol's strictness may have caused
under-counting of genuine but imprecisely stated findings; claims that could not
be assessed as fully specified were treated equivalently to refuted ones.
Third, the review covers only publicly accessible papers; internal safety
research at major AI laboratories is not represented.
Fourth, the rapidly evolving pace of the field means that findings reported here
may be superseded by developments after the review cutoff date of mid-2025.
Fifth, the absence of verified findings about audio and video safety reflects
the state of the literature, not a deliberate exclusion; systematic search was
conducted for these modalities.
Sixth, the inclusion boundary between reviewed systems and contrast baselines was applied using the following principle: a paper is included as a reviewed system if its primary contribution is a guardrail mechanism for the production inference pipeline (regardless of whether the mechanism operates at training or inference time); it is treated as a contrast baseline if its primary contribution is a training procedure whose safety benefit is incidental to another objective (e.g., fine-tuning for capability).
DRESS~\cite{chen2024dresssafety} is classified as a contrast baseline because its contrastive decoding approach was proposed as a generation quality improvement with safety as a secondary benefit; VLGuard~\cite{li2024vlguard} is classified as a reviewed system because its training objective is explicitly safety-oriented.
This principle was applied consistently across all included papers; a complete inclusion/exclusion record is available in Appendix~\ref{sec:appendix}.
